% !TeX root = ../main.tex

\chapter{Methodology}\label{ch:methodology}

% Scope rule for this chapter: state what was built and how it would be judged.
% Every outcome, verdict and void rung belongs in Chapter 6.

\section{Pre-Existing Framework, Research Questions and Scope}\label{sec:preexisting}

This thesis was conducted within an ongoing research programme on individualized
functional-connectivity abnormality detection. Prior to the start of the thesis, an
existing methodological and software framework for \ac{fMRI} preprocessing,
functional-connectivity analysis, subject-level graph construction, and
graph-autoencoder-based modelling was available. The present thesis builds on this
foundation and applies and extends it toward longitudinal \ac{MCI}-to-\ac{AD} conversion
prediction, cross-cohort evaluation, and the longitudinal modelling and benchmarking
experiments described below.

This chapter states what was built and how it would be judged.
\autoref{sec:problem-formulation} defines the formal prediction task and the graph
constructed on top of the connectivity matrix \autoref{ch:data} describes.
\autoref{sec:static-representation} covers static representation learning via the
\ac{GAAE} and \ac{VGAE}. \autoref{sec:spatial-first} describes the spatial-first
trajectory models that build on those representations. \autoref{sec:tfgn} introduces the
temporal-first alternative and its design rationale. \autoref{sec:protocol} specifies the
experimental protocol that judges all arms. \autoref{sec:instrumentation} describes the
interpretability and robustness instrumentation.

The four research questions that organise the experimental programme are as follows.
\textbf{RQ1}: does longitudinal modelling add information beyond static
functional-connectivity representations, answered in \autoref{sec:rq1}?
\textbf{RQ2}: does spatial-first graph representation learning add value, answered in
\autoref{sec:rq2}?
\textbf{RQ3}: does temporal-first modelling improve over spatial-first modelling, answered
in \autoref{sec:rq3}?
\textbf{RQ4}: do the findings transfer to an independent cohort, answered in
\autoref{sec:rq4}?
RQ4 is framed as transfer rather than as validation: \ac{OASIS} is an independent
external cohort, not a held-out sample from the same distribution as the training data,
so its results inform cross-cohort generalisability rather than a classical validation claim.

\section{Problem Formulation and Graph Construction}\label{sec:problem-formulation}

Subject-level functional-connectivity graph modelling was already part of the pre-existing
project framework. The present analyses use this representation for the longitudinal
experiments described below. The staged spatial-first pipeline likewise formed part of the
pre-existing modelling framework. In this thesis, the consequences of this design are
treated as an explicit hypothesis and tested through controlled ablations and comparison
with the temporal-first alternative.

\paragraph{Formal statement.}
A subject $i$ contributes a variable-length sequence of $T_i$ visits. Each visit $t$
yields a graph $G_i^{(t)} = (V, E_i^{(t)}, X_i^{(t)})$, where the node set $V$ is fixed
at $|V| = 200$ regions of interest under the Schaefer-200 parcellation,
$E_i^{(t)} \subseteq V \times V$ is a binary edge set constructed from the visit's
functional-connectivity matrix (described below), and $X_i^{(t)} \in \mathbb{R}^{200
\times 200}$ is the matrix of node feature vectors. The prediction target is a single
binary subject-level label $y_i \in \{0, 1\}$, with $y_i = 1$ for a converter and
$y_i = 0$ for stable \ac{MCI}, as defined in \autoref{sec:delcode}. The effective sample
size for any downstream classifier is therefore the number of subjects, not the number of
scans, because the head produces one prediction per subject from the full visit sequence.
This constraint on the denominator is the primary motivation for the capacity arguments in
\autoref{sec:spatial-first} and \autoref{sec:tfgn}.

\paragraph{Node features.}
Each \ac{ROI} $j$ is represented by its own row from the $200 \times 200$
Fisher-$z$-transformed functional-connectivity matrix that \autoref{sec:preprocessing}
describes. The node-feature matrix $X_i^{(t)} \in \mathbb{R}^{200 \times 200}$ is
therefore the full connectivity matrix, so that each node carries the profile of its
pairwise coupling with every other region. The input feature dimensionality is consequently
$\texttt{in\_features} = 200$.

\paragraph{Edge construction.}
The binary adjacency matrix $A_i^{(t)}$ is constructed by symmetrised $k$-nearest-neighbour
sparsification over the absolute values of the functional-connectivity matrix. For each
\ac{ROI} $j$, the $k$ largest absolute entries of its connectivity row are retained and
the diagonal is excluded. The result is then symmetrised as $A \leftarrow \max(A, A^\top)$,
so that if \ac{ROI} $j$ is among \ac{ROI} $l$'s top-$k$ neighbours then the edge $(j, l)$
is present regardless of whether the converse holds. Taking absolute values retains strongly
anti-correlated pairs as connected nodes, which avoids discarding the negatively coupled
region pairs that are known to be functionally meaningful in resting-state data. Edge
attributes are binary. The value of $k$ differs by stage: $k = 16$ during \ac{GAAE} and
\ac{VGAE} pretraining, and $k = 8$ for \ac{GEC}, \ac{GEP}, \ac{GELSTM}, \ac{GEGRU} and
\ac{TFGN} supervised training.

\paragraph{Time conditioning.}
The three cohorts encode visit time differently: \ac{DELCODE} stores protocol months along
a nominal six-month-to-annual grid, while \ac{ADNI} and \ac{OASIS} store elapsed days from
each subject's baseline session. Both representations are converted to a common inter-visit
interval in months. The first visit always receives $\Delta t_1 = 0$, and subsequent
intervals are computed as
\[
  \Delta t_t = \frac{\text{cum}_t - \text{cum}_{t-1}}{108},
\]
where $\text{cum}_t$ is the cumulative months elapsed since baseline and the normaliser
(\texttt{MAX\_INTERVAL\_MONTHS} $= 108.0$) corresponds to nine years, the longest observed
follow-up in the pooled cohort. This representation directly addresses the irregular
longitudinal sampling challenge identified in \autoref{sec:sampling}: the model receives
the actual inter-visit gap rather than assuming a fixed cadence. The \ac{TFGN} instead
uses cumulative $\log(1 + \text{months})$ standardised with training-fold statistics,
because the logarithmic compression reduces the dynamic range of the temporal input without
losing the ordering information.

\paragraph{Worked example.}
Consider a subject with three visits at 0, 12 and 30 months from baseline. The first visit
produces a graph $G^{(1)}$ with $\Delta t_1 = 0$, a $200 \times 200$ node-feature matrix
$X^{(1)}$ drawn from that session's Fisher-$z$ connectivity matrix, and a binary adjacency
$A^{(1)}$ retaining the eight largest absolute connections per \ac{ROI}. The second visit
yields $\Delta t_2 = (12 - 0)/108 \approx 0.111$, with its own $X^{(2)}$ and $A^{(2)}$.
The third visit yields $\Delta t_3 = (30 - 12)/108 \approx 0.167$. The recurrent models
receive the sequence $\bigl[(X^{(1)}, A^{(1)}, 0),\, (X^{(2)}, A^{(2)}, 0.111),\,
(X^{(3)}, A^{(3)}, 0.167)\bigr]$ and produce a single binary prediction from the terminal
hidden state.

\section{Static Representation Learning: GAAE and VGAE}\label{sec:static-representation}

The \ac{GAAE} stage builds on a graph-attention-autoencoder implementation available in
the pre-existing project codebase. The configuration and modifications used for the
experiments reported in this thesis are described below. For the pooled \ac{GAAE}, the
change concerns the training cohort and experimental setting rather than the
graph-autoencoder concept itself.

\paragraph{GAAE architecture.}
The \ac{GAAE} encoder consists of three GATv2Conv layers \parencite{brody2022gatv2} with
two-dimensional edge features and residual connections, reducing the 200-dimensional input
to a 64-dimensional latent vector. The three layers have hidden dimensionalities
$200 \to 128 \to 128$ with two attention heads each, followed by batch normalisation,
ReLU activation and dropout at probability 0.3. The latent vector is then modulated by a
Feature-wise Linear Modulation (FiLM) layer conditioned on a two-dimensional covariate
vector of normalised age and binary sex encoding, so the encoder can adjust its per-dimension
scale and shift to account for demographic variation without requiring separate model
instances. A mirrored three-layer GATv2 decoder reconstructs the node-feature matrix, and a
separate inner-product adjacency decoder reconstructs the edge structure as
$\hat{A} = \sigma(ZZ^\top)$, where $Z \in \mathbb{R}^{200 \times 64}$ is the
node-level latent matrix and $\sigma$ denotes the logistic function applied element-wise.
The training loss is
\begin{equation}
  \mathcal{L}_{\text{GAAE}}
  = \operatorname{MSE}(\hat{X}, X)
  + \lambda_{\text{adj}}\,\operatorname{BCE}(\hat{A}, A),
  \label{eq:gaae-loss}
\end{equation}
where $\hat{X}$ is the reconstructed node-feature matrix, $X$ is the input, $\hat{A}$ is
the reconstructed adjacency, $A$ is the true binary adjacency, $\operatorname{MSE}$ is the
mean squared error over node features, $\operatorname{BCE}$ is the binary cross-entropy
over adjacency entries, and $\lambda_{\text{adj}} = 0.2$ weights the adjacency term. The
loss is masked so that block-diagonal batch padding does not contribute. Key
hyperparameters: latent
dimensionality 64, hidden dimensionality 128, two attention heads, dropout 0.3, learning
rate $10^{-3}$, weight decay $10^{-3}$, up to 500 training epochs, early-stopping patience
25 on validation loss, and $k = 16$ for the $k$-nearest-neighbour graph. The encoder has
952,136 parameters.

\paragraph{The pivot from reconstruction to classification.}
Because within-cohort between-subject variation in functional connectivity can exceed
between-group variation at the available sample sizes, as \autoref{sec:statistics}
quantifies, a reconstruction objective is not guaranteed to allocate encoder capacity to
the group contrast. The pipeline therefore separates pretext learning from the group
decision: the encoder is pretrained to reconstruct graph structure and then frozen, and a
separate classification head is trained on the resulting embeddings. The reconstruction
phase is treated as a representation-learning step whose value is tested empirically rather
than assumed; the results are reported in \autoref{sec:rq2}.

\paragraph{GEC and GEP as static classification heads.}
The \ac{GEC} adapter implements embed-then-classify: the frozen \ac{GAAE} encoder is
applied to a single visit's graph, the 200 node embeddings are mean-pooled to a
64-dimensional graph embedding, and a multilayer perceptron maps this embedding to a
classification logit. An attention-pooling variant (\texttt{GraphEncoderClassifierAttention})
replaces mean pooling with a learned gate network. The \ac{GEP} adapter is a multi-visit extension: it encodes each of a subject's visits
independently and mean-pools the resulting sequence of per-visit graph embeddings to a
single subject-level vector before classification, so it is a static model that uses visit
averaging rather than recurrence.

\paragraph{VGAE architecture and the free-bits objective.}
The \ac{VGAE} shares the GATv2 backbone of the \ac{GAAE} encoder but replaces the
deterministic output with two parallel convolution heads emitting a mean vector
$\mu \in \mathbb{R}^{64}$ and a log-variance vector $\log\sigma^2 \in \mathbb{R}^{64}$.
During training, a latent sample is drawn by the reparameterisation trick,
$z = \mu + \sigma \odot \varepsilon$ with $\varepsilon \sim \mathcal{N}(0, I)$; during
evaluation the mean $\mu$ is used directly. The FiLM conditioning is applied to $\mu$
after sampling, so that the conditioning shift is not penalised by the Kullback--Leibler
divergence term. An
optional feature decoder adds an auxiliary node-feature reconstruction head alongside the
inner-product adjacency decoder. The standard Kullback--Leibler term can drive individual
latent dimensions to collapse to the prior, producing a representation where those
dimensions carry no information. The \ac{VGAE} therefore uses a free-bits objective: for
each dimension $d$ independently,
\begin{equation}
  \mathcal{L}_{\text{KL},d}
  = \max(\operatorname{KL}_d,\, b)
    + \max\!\bigl(b - \operatorname{KL}_d,\, 0\bigr)^2,
  \label{eq:free-bits}
\end{equation}
where $\operatorname{KL}_d$ is the per-dimension Kullback--Leibler divergence and $b$ is a
floor threshold. At or above the floor, \autoref{eq:free-bits} reduces to the ordinary
Kullback--Leibler term and has the same gradient. Below the floor, the first term is the
constant $b$ with zero gradient; the second term is positive with a negative gradient with
respect to $\operatorname{KL}_d$, which pushes the divergence back toward the floor rather
than allowing the dimension to remain collapsed. A plain clamp alternative sets the loss to
$b$ below the floor with zero gradient and therefore cannot pull a collapsed dimension
back into use; a softplus relaxation is also implemented and its behaviour is described in
\autoref{ch:app-implementation}. The anti-collapse hyperparameters are as follows: free-bits floor 1.23 and 1.44, $\beta$ weighting 0.00111 and 0.0005,
warmup over 41 and 43 epochs, optional feature-loss weight 2.93 and 0.5 for the two
final configurations.

\paragraph{Reconstruction fidelity as a method-level metric.}
We define reconstruction fidelity as the Pearson correlation coefficient between the
flattened input matrix and its flattened reconstruction, computed by a single canonical
function. Given input
$X$ and reconstruction $\hat{X}$, both are flattened to one-dimensional vectors
$\mathbf{a}$ and $\mathbf{b}$, and the correlation is
\begin{equation}
  r = \frac{\sum_i (a_i - \bar{a})(b_i - \bar{b})}
           {\sqrt{\sum_i (a_i - \bar{a})^2}
            \cdot \sqrt{\sum_i (b_i - \bar{b})^2}},
  \label{eq:pearson-r}
\end{equation}
where $\bar{a}$ and $\bar{b}$ are the means of the flattened vectors and the sum runs over
every matrix element. We use correlation rather than mean squared error because the
Fisher-$z$ connectivity scale is dataset-specific, so a scale-free correlation is
comparable across cohorts where an absolute reconstruction error is not. Fidelity is
reported as a cohort median with interquartile range computed over per-subject values. The
reconstruction targets differ by model: for the \ac{GAAE}, $X$ is the continuous
node-feature matrix and $\hat{X}$ is the feature decoder's output; for the \ac{VGAE},
$X$ is the binary adjacency matrix and $\hat{X}$ is the continuous edge-probability output
$\sigma(ZZ^\top)$. These two correlation values measure different reconstruction targets
and are not interchangeable as statements about which model reconstructs connectivity
better. The values are reported in \autoref{sec:rq2}.

\section{Spatial-First Trajectory Modelling}\label{sec:spatial-first}

This part builds on the pre-existing spatial-first graph-learning and graph-autoencoder
framework available at the start of the thesis and is applied here to longitudinal \ac{MCI}
conversion prediction. The encoder handling inherited from the pre-existing implementation
was modified during the thesis to enable the controlled fine-tuning ablation described
below.

\paragraph{The recurrent core.}
At each visit $t$, the shared encoder produces a visit embedding $z_t \in \mathbb{R}^{d}$,
where $d = 64$ when an encoder is present and $d = 200$ when the \texttt{none} arm
bypasses the encoder entirely and mean-pools the raw node features. The time-augmented
embedding $[z_t \,\|\, \Delta t_t]$ is consumed by a single-layer \ac{LSTM} with hidden
dimensionality 32, which processes the visit sequence and reads out the terminal hidden
state $h \in \mathbb{R}^{32}$. A classifier head of the form
$\text{Linear}(32) \to \text{ReLU} \to \text{Dropout}(0.3) \to \text{Linear}(1)$
maps $h$ to a classification logit. LayerNorm is applied before the head, so
variable-length evaluation at batch size one is numerically stable. The \ac{GEGRU} uses the
same architecture; a gated recurrent unit has three gates
against the \ac{LSTM}'s four, which gives approximately 25\% fewer recurrent parameters at
equal hidden dimensionality.

\paragraph{Regularisation and normalisation.}
Per-visit embeddings are z-scored using persistent mean and variance buffers fitted on the
training fold only, which is the recurrent analogue of the per-fold
\texttt{StandardScaler} used by the multilayer perceptron classification heads. Additional
regularisation includes dropout at probability 0.3, gradient clipping at a global norm of
1.0, weight decay $10^{-4}$, and cost-sensitive class weighting to handle the converter
minority, as described in \autoref{sec:protocol}. Early stopping terminates training when
validation loss fails to improve for 15 consecutive epochs.

\paragraph{Capacity argument.}
The effective sample size for the downstream heads is the number of subjects, not the
number of scans, because a single binary label is fitted per subject. In the \ac{DELCODE}
cross-validation pool there are 133 subjects, with approximately 100 in each training
fold. A hidden-128, two-layer \ac{LSTM} configuration accumulates 240,257 trainable
parameters in the recurrent and classifier components, corresponding to approximately 1,806
parameters per cross-validation subject. The breakdown is: \ac{LSTM} layer 1,
$4 \times 128 \times (65 + 128)$ plus biases $= 99{,}840$ parameters; layer 2,
$4 \times 128 \times (128 + 128)$ plus biases $= 132{,}096$ parameters; classifier head
$\approx 8{,}321$ parameters. The \ac{LSTM} component alone exceeds the entire multilayer
perceptron baseline in parameter count. The deployed configuration uses hidden
dimensionality 32 and one layer, reducing the recurrent-plus-head parameter count to
approximately 13,761 (roughly 103 per cross-validation subject). A \ac{GEGRU} variant at
the same hidden size reduces recurrent parameters by approximately 25\% further. These
figures are the design
rationale for the chosen specification: the recurrent inductive bias, not parameter count
alone, is the leading overfitting risk at this sample size, because a carried hidden state
applied at every visit and read out at the final step is exactly the function class that
can memorise a small cohort.

\paragraph{The four encoder arms and the RQ2 ablation.}
The value of reconstruction pretraining is tested by a single configuration knob,
which selects one of four arms:

\begin{center}
\small
\setlength{\tabcolsep}{4.5pt}
\begin{tabular}{lcccp{0.36\linewidth}}
\toprule
Arm & Built? & Loaded? & Trained? & Isolates \\
\midrule
\texttt{pretrained\_frozen} & yes & yes & no & Reference (pre-existing model) \\
\texttt{pretrained\_finetuned} & yes & yes & yes & Value of task adaptation \\
\texttt{random} & yes & no & yes & Value of reconstruction pretraining \\
\texttt{none} & no & no & --- & Value of the graph encoder itself \\
\bottomrule
\end{tabular}
\end{center}

The arms form a chain of contrasts: \texttt{none} versus \texttt{random} isolates the
encoder architecture, \texttt{random} versus \texttt{pretrained\_frozen} isolates
reconstruction pretraining, and \texttt{pretrained\_frozen} versus
\texttt{pretrained\_finetuned} isolates task adaptation. Under the \texttt{none} arm, no
encoder module is constructed; \texttt{encode\_visit} mean-pools the raw node features
instead, so the per-visit embedding width becomes 200 and the recurrent input width is
derived rather than hardcoded. These four arms apply only to the \ac{GELSTM} and
\ac{GEGRU} adapters, because the \ac{GEC} and \ac{GEP} adapters pre-compute embeddings
offline with a frozen encoder; end-to-end encoder training has no meaning for them without
restructuring those adapters.

\paragraph{Pre-specified interpretation table.}
The following interpretation table was written before any arm was run, so that the reading
of the results could not be fitted retrospectively. 

\begin{center}
\small
\begin{tabular}{p{0.38\linewidth}p{0.27\linewidth}p{0.27\linewidth}}
\toprule
Observation & Meaning & Consequence \\
\midrule
\texttt{random} $\approx$ \texttt{pretrained\_frozen} $\approx$ \texttt{pretrained\_finetuned} &
  Reconstruction adds nothing beyond the architecture prior &
  Keep as a compute-saving initialisation; report as a null finding \\
\texttt{none} $\approx$ every other arm &
  Graph encoder not earning its place &
  Focus should be about the longitudinal head; \\
\texttt{pretrained\_*} $>$ \texttt{random}, both clearly above \texttt{none} &
  Pretraining contributes real transferable structure &
  Keep the \ac{GAAE} stage; quantify the gain \\
\texttt{pretrained\_finetuned} $>$ \texttt{pretrained\_frozen} &
  Pretrained features under-adapted to conversion task &
  Switch to finetuned arm; monitor overfitting \\
\texttt{pretrained\_finetuned} $<$ \texttt{pretrained\_frozen} &
  Fine-tuning overfits the small converter cohort &
  Justify freezing explicitly as a regularisation choice \\
\texttt{random} $>$ \texttt{pretrained\_*} &
  Reconstruction pretraining actively harmful &
  Investigate pretext--downstream mismatch \\
All four arms $\approx$ metadata floor &
  Imaging not contributing signal at this sample size &
  Escalate to sample size and label definition \\
\bottomrule
\end{tabular}
\end{center}

The results of applying this table are in \autoref{sec:rq2}. The scope limit of this
ablation is that it applies only to the \ac{GELSTM} and \ac{GEGRU} adapters; the
equivalent question for \ac{GEC} and \ac{GEP} is a separate follow-up.

\section{Temporal-First Modelling: The TFGN Architecture}\label{sec:tfgn}

\paragraph{Hypothesis.}
Every longitudinal classifier described in \autoref{sec:spatial-first} pools a visit's
whole-brain graph to a single 64-dimensional vector before any temporal model sees it.
Region identity is consequently discarded before any trajectory is modelled: the recurrent
core operates on a compressed representation that mixes all regional signals together.
\ac{TFGN} defers this pooling. A node-shared recurrent core first encodes each of the 200
regions' own connectivity-row trajectory independently, and pooling is applied only
afterward, optionally after a graph propagation stage. If region-resolved temporal dynamics
carry conversion-relevant signal that early pooling destroys, deferring pooling should
improve prediction at matched or lower capacity. 

\paragraph{Inputs.}
The per-subject input to \ac{TFGN} consists of the visit stack
$X \in \mathbb{R}^{T \times 200 \times 200}$, the cumulative log inter-visit interval
$\log(1 + \text{months})$ standardised with training-fold statistics, the baseline
adjacency $A_0$ constructed as a binary $k$-nearest-neighbour graph with $k = 8$ on the
absolute baseline connectivity matrix, a covariate vector $[\text{age},\, \text{sex}]$,
and, for arms that use them, the per-subject binary connectivity-change mask $M$, the
strength centrality vector of $|A_0|$, and the rank-sigmoid drift anchor $\tilde{d}$.

\paragraph{Switchable stages.}
The architecture is implemented as a sequence of seven independently switchable modules,
controlled by a single \texttt{TFGNTrainConfig} object, so that the entire ablation ladder
corresponds to toggling knobs on one model.

\textbf{(1) Node-shared LSTM.} A single \ac{LSTM} with hidden dimensionality 64, one
layer and dropout 0.3, applied with shared weights to each of the 200 nodes' own
connectivity-row sequence. Each node $i$ has its own trajectory of 200-dimensional
feature vectors (the rows of $X^{(t)}$), and the shared \ac{LSTM} produces a
node-specific terminal hidden state $h_i^{(T)} \in \mathbb{R}^{64}$. The log
inter-visit interval is optionally concatenated to the node input at each step.

\textbf{(2) Temporal saliency gate.} A per-node scalar gate
\[
  s_i = \sigma\!\bigl(w^\top \operatorname{LeakyReLU}(W_s h_i)\bigr)
\]
is applied as residual scaling $(1 + s_i)\,h_i$, where $W_s \in \mathbb{R}^{d_s \times
64}$ and $w \in \mathbb{R}^{d_s}$ are learned parameters and $\sigma$ is the logistic
function. A gate value of zero passes the hidden state unchanged; values above zero amplify
the contribution of regions whose temporal trajectory is more informative. The gate is
regularised by a sparsity penalty described in the loss paragraph below.

\textbf{(3) GVAE encoder.} When a reconstruction target is set, a \ac{GVAE} is constructed with GATv2
heads for the mean $\mu$ and log-variance $\log\sigma^2$, hidden dimensionality 128,
latent dimensionality 64 and two attention heads, operating over $A_0$. FiLM conditioning
of $\mu$ on $[\text{age},\, \text{sex}]$ is applied before reparameterised sampling. Under
\texttt{recon\_target: none} this stage is absent entirely.

\textbf{(4) Fusion.} When a \ac{GVAE} is present, the temporal hidden state $h_i$ and
the graph latent $z_i$ are combined by either
$\operatorname{LayerNorm}(W_u[h_i \,\|\, z_i])$ (\texttt{concat\_residual}) or by
passing $z_i$ alone (\texttt{z\_only}).

\textbf{(5) Readout.} The 200 node representations are reduced to a single subject
embedding by mean pooling or by attentive pooling using exact sparsemax
\parencite{martins2016sparsemax}. Sparsemax produces true zeros rather than a dense soft
distribution, so nodes assigned zero weight contribute nothing to the subject embedding.

\textbf{(6) Dual-score head.} The classification head emits both a binary classification
logit and a per-node topological saliency score $s_{\text{topo}}$, anchored to the
strength centrality of $|A_0|$ by an auxiliary mean-squared-error term weighted by
$\lambda_{\text{cent}}$. This stage is the interpretability layer; see
\autoref{sec:instrumentation}.

\textbf{(7) Cohort-adversary head.} Used only by the escalation arm. A gradient-reversal
layer, which applies the identity in the forward pass and negates and scales gradients in
the backward pass, connects the pooled subject embedding to a binary
ADNI-versus-DELCODE classification head, training the encoder to be uninformative about
cohort identity.

\paragraph{Composite loss.}
The training objective is
\begin{equation}
  \mathcal{L}
  = \mathcal{L}_{\text{BCE}}
  + \lambda_{\text{sparse}}\,\operatorname{KL}(\bar{s} \,\|\, \rho)
  + \lambda_{\text{drift}}\,\operatorname{MSE}(\bar{s},\, \tilde{d})
  + \lambda_{\text{cent}}\,\operatorname{MSE}(s_{\text{topo}},\, c)
  + \mathcal{L}_{\text{recon}}
  + \beta\,\mathcal{L}_{\text{KL}},
  \label{eq:tfgn-loss}
\end{equation}
where $\mathcal{L}_{\text{BCE}}$ is the cost-weighted binary cross-entropy on the
classification logit; $\operatorname{KL}(\bar{s} \,\|\, \rho)$ is the
Kullback--Leibler divergence of the mean gate activation $\bar{s}$ from the target
sparsity $\rho = 0.15$, which encourages at most 15\% of nodes to have above-baseline
activations; $\operatorname{MSE}(\bar{s},\, \tilde{d})$ anchors the mean gate activation
to the rank-sigmoid drift anchor $\tilde{d}$; $\operatorname{MSE}(s_{\text{topo}},\, c)$
anchors the topological saliency score to the strength centrality $c$ of $|A_0|$;
$\mathcal{L}_{\text{recon}}$ is the reconstruction loss on the change mask (when active);
and $\beta\,\mathcal{L}_{\text{KL}}$ is the warmed-up \ac{GVAE} Kullback--Leibler term,
computed on pre-FiLM $\mu$ using the same free-bits function as \autoref{eq:free-bits} so
that the conditioning shift is not penalised by the prior. The drift anchor weight is set
as $\lambda_{\text{drift}} = 0.1\,\lambda_{\text{sparse}}$, so the sparsity prior
dominates at the margin.

\paragraph{Pre-specified design decisions.}
The following decisions were fixed prior to any run.

\textbf{Reconstruction target.} The inner-product output $\sigma(ZZ^\top) \in (0,1)$
cannot be trained with binary cross-entropy against a raw connectivity-change matrix
$\Delta A \in [-2,2]$. The pre-specified default reconstruction target is
\texttt{delta\_a\_topk}: a per-subject binary change mask $M_{ij} = \mathbb{1}[|\Delta
A_{ij}| \geq q_{1-\kappa}]$ with $\kappa = 0.10$, so the top 10\% of changed edges are
positive. The loss uses binary cross-entropy with a positive-edge weight of
$(1 - \kappa)/\kappa = 9$ to compensate for the strong class imbalance. Construction
raises an error rather than proceeding silently when any subject's positive-edge density
falls outside $[0.01,\, 0.5]$.

\textbf{Topological anchor.} Eigenvector centrality is ill-defined on signed connectivity
matrices because negative weights violate the Perron--Frobenius conditions that guarantee
a unique dominant eigenvector. The anchor is therefore the strength centrality of $|A_0|$,
computed as the column sum of the absolute-value adjacency and z-scored with training-fold
statistics.

\textbf{Drift anchor.} A raw drift magnitude $\|x_i^{(T)} - x_i^{(1)}\|_2$ is a dense
positive scalar that would conflict with the sparsity Kullback--Leibler term. The anchor
is instead a rank-sigmoid transformation of the within-subject drift:
\begin{equation}
  q_i = \frac{\operatorname{rank}(d_i)}{N - 1},
  \quad
  \tilde{d}_i = \sigma\!\left(\frac{q_i - (1 - \rho)}{\tau_d}\right),
  \quad \tau_d = 0.05,
  \label{eq:drift-anchor}
\end{equation}
where $d_i = \|x_i^{(T)} - x_i^{(1)}\|_2$, $N$ is the number of subjects in the
current batch, $\rho = 0.15$ is the target sparsity, and $\sigma$ is the logistic
function. By construction $\operatorname{mean}(\tilde{d}) \approx \rho$, so the drift
anchor and the sparsity Kullback--Leibler target agree on the expected gate activation
level. The rank transformation also makes the anchor scale-free across cohorts with
different connectivity amplitudes.

\textbf{Strict determinism.} The GATv2 scatter-based aggregation and the sparsemax
backward pass have non-deterministic GPU kernels. Without forcing determinism, the
seed-level standard error of cross-validation results would partially measure GPU
scheduling noise rather than model behaviour.

\paragraph{Cohort-shift control as a design decision.}
Cohort identity is deliberately withheld from the model: a FiLM cohort covariate would
make \ac{OASIS} an untrained one-hot category at inference time and its behaviour would
be unspecified. Instead, every pooled run computes a mandatory cohort probe: a logistic
regression decodes cohort identity from the out-of-fold subject latents, and its \ac{AUC}
is recorded. If this probe exceeds 0.75 on the selected
arm, the pre-specified response is to re-run that arm with the gradient-reversal
cohort-adversary head and report both alongside each other. The outcome of this check is
reported in \autoref{sec:rq4}.

\paragraph{Self-supervised pretraining.}
Two pretraining runs are conducted, both on pooled \ac{ADNI} plus \ac{DELCODE} only;
\ac{OASIS} is excluded from everything upstream so that it remains genuinely external.

The first run (P1) retrains the \ac{GAAE} on the pooled unlabelled graph collection
(approximately 3,700 scans), using an architecture identical to the existing
\ac{DELCODE}-only checkpoint (hidden 128, latent 64, two attention heads, $k = 16$,
adjacency-loss weight 0.2, 500 epochs). Keeping the architecture identical ensures that
downstream arms using this checkpoint remain comparable to arms using the \ac{DELCODE}-only
checkpoint.

The second run (P2) trains the node-shared \ac{LSTM} on a self-supervised next-visit
forecasting task: given each of the 200 nodes' connectivity-row history up to visit $t$,
the model predicts the connectivity row $x_i^{(t+1)}$ at the next visit, trained with mean
squared error loss over every subject with at least two visits in the pooled pretraining
split. A persistence baseline, predicting $x^{(t+1)} = x^{(t)}$, is defined as the method
comparison; its mean squared error is reported alongside the trained model's in
\autoref{sec:rq3}.

\paragraph{Capacity accounting.}
The same subject-level denominator applies here. In the pooled cross-validation pool there
are 248 subjects. The \ac{TFGN} node-shared \ac{LSTM} at hidden 64 and one layer has the
recurrent parameter footprint of a single \ac{LSTM} at those dimensions, since weights are
shared across all 200 nodes. The architecture is kept at a capacity level comparable to
the spatial-first models at hidden 32, so that differences in performance can be attributed
to the pooling-order inversion rather than to capacity differences. No parameter-count
comparison that identifies a configuration as superior is made here; those comparisons are
in \autoref{sec:rq3}.

\section{Experimental Protocol}\label{sec:protocol}

\paragraph{Class imbalance and threshold policy.}
Converters are a minority of the cohort in all three datasets. We address this by
cost-sensitive class weighting of the binary cross-entropy loss: class weights are set to
the inverse class frequencies in the training fold and recomputed per fold. Resampling
methods are not used because they add no information about the minority class at these
sample sizes. The decision threshold is selected as the value maximising F1 on out-of-fold
predictions, stored in the model checkpoint, and reused on the test set without any
re-thresholding. One point on nomenclature: the legacy configuration dataclass defaults the threshold mode to Youden's J,
but every runner configuration and every experiment reported in this thesis uses the
best-F1 mode. The best-F1 selection is the protocol default.

\paragraph{Cross-validation.}
All supervised experiments use five-fold stratified group cross-validation, stratified on
the binary label and grouped on subject identifier so that no subject's visits appear in
more than one fold. The per-fold \texttt{StandardScaler} for embeddings is fitted on the
training fold only and applied to validation and test embeddings. The log inter-visit
interval and strength centrality statistics for \ac{TFGN} arms are fitted on the training
fold and carried inside the saved model state. All experiments use seeds 42, 43, 44 and 45
to support seed-level variance estimation.

\paragraph{Training.}
The optimiser is Adam with learning rate $10^{-3}$, batch size 16, and gradient clipping
at a global norm of 1.0. Training runs for up to 100 epochs with early stopping at patience
20. A \texttt{ReduceLROnPlateau} scheduler halves the learning rate when validation loss
fails to improve for five consecutive epochs, down to a floor of $10^{-6}$. Cost-sensitive
class weighting applies to the binary cross-entropy throughout.

\paragraph{Ablation ladder as a design.}
The ablation ladder is a sequential architecture search in which each rung inherits the
kept configuration of all preceding rungs and changes exactly one knob from the last
surviving state. A rung that fails the keep-or-drop criterion has its knob reverted and
the next rung branches from the last surviving configuration. The two design contrasts of
primary interest are stated in advance: S0c against S1, in which neither encoder is
pretrained, isolates the pooling-order inversion; S0b against S1c, in which both encoders
are self-supervised, isolates the inversion against a matched pretraining baseline. The
ladder arms are as follows.

\textbf{S0a} is a logistic regression over principal components of the flattened
connectivity-change matrix concatenated with visit count, total follow-up months, age and
sex. It is a linear floor testing whether signal exists in raw connectivity change.

\textbf{S0b} is the spatial-first \ac{GELSTM} with the pooled \ac{GAAE} checkpoint frozen.
It is the matched comparator for S1c: both S0b and S1c use a self-supervised encoder.

\textbf{S0c} is the same spatial-first \ac{GELSTM} with random encoder initialisation,
trained end-to-end. It is the matched comparator for S1: neither S0c nor S1 uses
self-supervised pretraining.

\textbf{S0d} reproduces the Brain-TokenGT competitor arm \parencite{dong2023braintokengt}
as a reference point, characterised in \autoref{sec:braintokengt}. Its scatter-gather
operations do not engage deterministic algorithms, so
its same-seed test \ac{AUC} has been observed to vary substantially across runs; it is
reported with that caveat rather than as a controlled baseline.

\textbf{S1} uses a randomly initialised node-shared \ac{LSTM} with the gate disabled, no
reconstruction target, \texttt{fusion: z\_only} and mean readout. The S0c-against-S1
comparison isolates the architectural flip with no other confound.

\textbf{S1b} repeats S1 with the node-shared \ac{LSTM} initialised from the P2
self-supervised checkpoint. It tests whether node-level forecasting pretraining is
beneficial.

\textbf{S1c} adds a \ac{GVAE} reconstruction stage on top of S1's carried-forward configuration. The
S0b-against-S1c comparison tests the pooling-order inversion when both encoders are
self-supervised.

\textbf{S2} uses gate with $\lambda_{\text{sparse}} = 0.1$,
$\lambda_{\text{drift}} = 0.01$ and $\rho = 0.15$. It tests whether suppressing
temporally static regions is beneficial.

\textbf{S3} changes the fusion mode to \texttt{concat\_residual}. Under S1's configuration there is no \ac{GVAE} stage and therefore no
graph latent $z_i$ to fuse with the temporal hidden state $h_i$. The fusion knob
consequently has no gradient path to any new parameters under this configuration and the
rung is structurally inert, a property of the architecture under that knob setting.

\textbf{S4} replaces mean readout with sparsemax attentive pooling. It tests whether
learned node weighting improves over uniform pooling.

\textbf{S5} adds the dual-score head with $\lambda_{\text{cent}} = 0.1$. This rung is
pre-specified as kept regardless of its contribution to \ac{AUC}: a saliency map that does
not improve classification accuracy can still be a valid or invalid interpretability
artefact, and the question must not be dropped if the AUC result is neutral. The
interpretability validation is described in \autoref{sec:instrumentation}.

\textbf{SENS} repeats the selected arm at \texttt{min\_visits=3}, reducing the pool to
140 subjects. It tests whether longer sequences confer a larger advantage.

\textbf{W3} re-runs the \ac{GELSTM} and the \ac{TFGN} selected arm with
\texttt{max\_visits=3} to create a matched comparison with Brain-TokenGT, which is
architecturally constrained to at most three visits.

\paragraph{The four-tier evaluation protocol.}
This protocol was fixed in its entirety prior to reading any ladder result.

\textbf{Tier 1 (floor gates).} Every arm must exceed two pre-computed baselines: the
demographics-only logistic regression (age and sex only), and the static single-visit
classifier evaluated per fold via a \texttt{fold\_probe} hook inside the cross-validation
runner, so that neither baseline requires reading the held-out test set.

\textbf{Tier 2 (selection rule).} The independent unit is the seed ($n = 4$), not the
fold, because folds within a seed share a model class and are not independent draws. For
each seed, the mean paired per-fold change in \ac{AUC} is computed (rung $k$ minus its
parent, matched fold for fold). A rung is kept if and only if the mean of the four
seed-level means exceeds their standard error:
$\operatorname{mean}(\Delta) > \operatorname{SE}(\Delta)$.
The two possible outcomes are asymmetric by construction: a failure means the effect is
undetectable at this sample size, never that it is harmful; a pass means the rung is worth
carrying forward, never that the effect is statistically significant. When the fold-matched
and pooled formulations of the criterion disagree, the pre-specified one-standard-error
tie-breaker applies: among kept arms, the simplest configuration within one standard error
of the best is preferred. This increases the evaluation sample size from 64 (in-domain test) to 248 (cross-validation pool),
reducing the minimum detectable effect from approximately 0.08 to approximately 0.04
\ac{AUC} units, and removes the risk of selecting on the held-out test data.

\paragraph{Worked example.}
Consider a rung evaluated at four seeds, with seed-level mean $\Delta$\ac{AUC} values of
$+0.031$, $+0.008$, $+0.024$ and $-0.005$. The mean of the four values is
$(0.031 + 0.008 + 0.024 - 0.005)/4 = 0.0145$. The standard deviation is approximately
$0.0147$ and the standard error is $0.0147/\sqrt{4} \approx 0.0074$. Since
$0.0145 > 0.0074$, the rung passes and its knob is carried forward. Had the fourth value
been $-0.025$ instead, the mean would be $0.0095$ and the standard error $0.0209$, giving
a failure; the result would be stated as ``undetectable at this sample size,'' not as
evidence that the mechanism is harmful.

\textbf{Tier 3 (robustness vetoes).} After a rung passes Tier 2, five checks with
pre-specified thresholds are applied. A rung fails if any of the following conditions
holds: (i) per-cohort out-of-fold \ac{AUC} falls to or below the demographics-only floor
for any cohort; (ii) the standard deviation of the decision threshold across the 20
fold-seed combinations exceeds 0.15; (iii) the expected calibration error after temperature
scaling fitted on out-of-fold predictions exceeds 0.10; (iv) the Spearman correlation
between scan count and classification score within stable subjects exceeds 0.30; or (v)
the cohort probe \ac{AUC} exceeds 0.75.

\textbf{Tier 4 (held-out estimation).} The 64-subject in-domain test set and the
60-subject \ac{OASIS} set are read exactly once, on the selected arm and its pre-specified
secondary arms, at the out-of-fold-derived threshold with no retraining.
\texttt{common/frozen\_read.py} reconstructs each arm's adapter from its saved summary and
checkpoint and refuses to score a split whose result keys already exist unless an explicit
overwrite is provided, so the one-shot read cannot be spent a second time.

\paragraph{Scope table.}
Spending the one-shot held-out splits on every ladder rung would turn the external
evaluation into a selection set. The following table states which arms were eligible for
which tier.

\begin{center}
\small
\begin{tabular}{lccc}
\toprule
Arm & OOF & In-domain test & External (OASIS-3) \\
 & ($n=248$) & ($n=64$) & ($n=60$) \\
\midrule
S0-demo, S0a, S0b, S0c & yes & no & no \\
S0d Brain-TokenGT & yes & ad hoc, post hoc & ad hoc, post hoc \\
S1 (primary) & yes & yes & yes \\
S1b (secondary) & yes & yes & yes \\
S1c both versions & yes & no & no \\
S2, S3, S4 & yes & no & no \\
S5 (secondary, interpretability) & yes & yes & yes \\
SENS, W3 arms & yes & no & no \\
\bottomrule
\end{tabular}
\end{center}

\paragraph{Resolvability statement.}
At the in-domain test size of $n = 64$ and the external size of $n = 60$, the empirical standard error of the \ac{AUC} is $\text{SE} \approx 0.04\text{--}0.045$, producing a 95\% confidence interval half-width ($1.96 \times \text{SE}$) of approximately $0.08\text{--}0.09$. Consequently, \ac{AUC} differences below approximately $0.08$ fall entirely within the sampling noise margin and are not resolvable with bootstrap confidence intervals. Because sample error scales with $1/\sqrt{N}$, the four-fold larger out-of-fold pool ($n = 248$) cuts the standard error and confidence interval width in half, establishing a resolvable floor of approximately $0.04$. Every keep-or-drop verdict is therefore a screening heuristic at four seeds, stating whether an effect is detectable at this sample size, never a significance test and never a claim that an undetected effect is absent. Configuration knob values and run identifiers are tabulated in \autoref{ch:app-reproducibility}.

\section{Interpretability and Robustness Instrumentation}\label{sec:instrumentation}

This section describes the instrumentation used to assess the trained models. All
findings from applying this instrumentation are reported in the corresponding subsections
of \autoref{ch:results}.

\paragraph{The disease axis and per-scan disease score.}
We construct a linear direction in the encoder's 64-dimensional latent space that separates
converters from stable \ac{MCI} subjects. Every training-plus-validation scan is encoded
with the frozen encoder and the resulting node-level embeddings are mean-pooled to a
single graph-level vector $z_i \in \mathbb{R}^{64}$. These are z-standardised per
dimension to obtain $Z_{\text{std}} \in \mathbb{R}^{N \times 64}$. An L2-regularised
logistic regression is then fitted on $(Z_{\text{std}},\, y)$, where $y \in \{0,1\}^N$ is
the subject-level label vector and $N$ is the number of training scans. The weight vector
$\tilde{w} \in \mathbb{R}^{64}$ is the normal to the decision hyperplane. Its unit
normalisation $\hat{w} = \tilde{w}/\|\tilde{w}\|_2$ is the disease axis. The per-scan
disease score is
\begin{equation}
  s_i = z_i \cdot \hat{w},
  \label{eq:disease-score}
\end{equation}
where $z_i$ is the standardised embedding of scan $i$. A score $s_i > 0$ places scan $i$
on the converter side of the decision boundary $s = 0$; $s_i < 0$ places it on the
stable-\ac{MCI} side. The disease score compresses the classifier's decision function to a
single number per scan, comparable to a continuous biomarker readout. The leakage contract:
the axis is supervised feature selection computed on training data only and is never
re-fitted on test or external data.

\paragraph{The orthogonal-residual trajectory space.}
Subtracting the disease-axis component from each embedding gives an axis-orthogonal
residual $R_i = z_i - s_i\,\hat{w}$. Principal component analysis applied to the residual
matrix yields two components $(\mathrm{PC}_1, \mathrm{PC}_2)$ capturing the largest
sources of within-class variation. A subject's visit-to-visit trajectory can be drawn in
the three-dimensional space $(s, \mathrm{PC}_1, \mathrm{PC}_2)$ against the decision
boundary at $s = 0$, giving a geometrically grounded visualisation of longitudinal change
that separates the conversion-relevant direction from orthogonal within-class variation.

\paragraph{Latent steering.}
A probe embedding $z_0$ is traversed along the disease axis in steps:
$z(\alpha) = z_0 + \alpha\,\hat{w}$ for $\alpha \in \{-3\sigma, -2\sigma, -1\sigma, 0,
+1\sigma, +2\sigma, +3\sigma\}$, where $\sigma$ is the standard deviation of the
disease scores across the training set. Each steered point is decoded back to a
connectivity matrix via the \ac{GAAE} decoder, and the difference map
$\text{FC}(\alpha) - \text{FC}(0)$ is computed for each step. The purpose is
methodological: a disease axis that is purely a discriminative artefact with no grounding
in connectivity structure would produce structureless difference maps; a direction that
captures connectivity change relevant to conversion would produce anatomically structured
maps, which are inspected as a post-hoc plausibility check. No inference is drawn from
this experiment alone.

\paragraph{Fisher discriminant ratio.}
The Fisher discriminant ratio per latent dimension $j$ is
\begin{equation}
  \operatorname{FDR}_j
  = \frac{(\mu_{1j} - \mu_{0j})^2}
         {\sigma_{1j}^2 + \sigma_{0j}^2 + \varepsilon},
  \label{eq:fdr}
\end{equation}
where $\mu_{1j}$ and $\mu_{0j}$ are the class means and $\sigma_{1j}^2$ and $\sigma_{0j}^2$
are the class variances of dimension $j$ over the training subjects, and $\varepsilon$ is a
small constant for numerical stability. The same statistic is mapped back to regions by
computing it over the 200-dimensional node-level embeddings before pooling, yielding a
per-region discriminability score. This is supervised feature selection computed on
training data only; the resulting maps are descriptive and do not influence any
classification decision.

\paragraph{Latent visualisation.}
\ac{UMAP} projections of the 64-dimensional graph embeddings described above, with the
disease-axis decision boundary overlaid, assess whether the two group clusters are
geometrically separated in the reduced space.

\paragraph{Four perturbation methods.}
Four strategies perturb held-out samples to measure decision stability, each targeting a
different failure mode.

\textbf{Feature noise} adds Gaussian noise
$\boldsymbol{\epsilon} \sim \mathcal{N}(\mathbf{0},\, (\sigma_X \cdot \lambda)^2 I)$
to the node-feature matrix $X$, where $\sigma_X$ is the empirical standard deviation over
all elements of $X$ and $\lambda \in [0,1]$ is the noise level. The graph topology is
held fixed. This tests sensitivity to measurement noise in the connectivity features.

\textbf{Matrix noise with graph rebuild} applies the same Gaussian perturbation to $X$ and
then recomputes the $k$-nearest-neighbour adjacency from the corrupted features, so the
graph topology changes as a consequence of feature corruption. This is the most
structurally disruptive strategy, simulating the scenario where noisy measurements produce
a different connectivity graph.

\textbf{Edge perturbation} leaves the node features unchanged and modifies only the graph
topology. A fraction $\lambda$ of existing edges is dropped and an equal number of random
$(u,v)$ pairs is added, with self-loops excluded. The retained and added edges are
deduplicated via \texttt{torch.unique} to prevent duplicate columns from violating the
binary adjacency assumption. This tests stability to connectivity graph scaffold changes
without any feature modification.

\textbf{Conditioning noise} perturbs the covariate vector: noise
$\delta_a \sim \mathcal{N}(0,\, (0.05\,\lambda)^2)$ is added to the normalised age $a$,
and $\delta_s \sim \mathcal{N}(0,\, (0.10\,\lambda)^2)$ is added to the encoded sex $s$,
with sex clipped to $[0,1]$ after perturbation. The graph and features are held fixed.
The scale factors (0.05 for age, 0.10 for sex) are smaller than the feature-noise scale
because these are normalised scalar conditioning inputs. This method tests whether the
decision is driven by the graph representation or by demographic covariates.

The sweep applies noise levels 0, 5, 10, 20 and 30\% with ten independent trials per
(subject, method, level) combination. The edge-perturbation deduplication fix and the
tests pinning it are described in \autoref{ch:app-implementation}.

\paragraph{Cohort stability rate.}
The reported robustness metric is the fraction of perturbation trials in which a subject's
classification decision survives intact, scored against a per-cohort one-versus-rest
threshold derived on the validation fold. Subjects are selected by their margin beyond that
threshold, so the test asks whether the most confident predictions are also the most stable
under perturbation. Mean error drift across perturbation trials is reported alongside the
stability rate.

\paragraph{Scanner-drift simulation.}
We simulate cross-cohort acquisition drift by perturbing \ac{DELCODE} test trajectories at
calibrated noise levels. This is stated as a simulation and not as a validation.

The calibration anchor is the standard deviation of edge-wise functional-connectivity
differences between consecutive same-subject \ac{DELCODE} visits, which represents the
level of variation the model already tolerates from normal longitudinal change. Each target
cohort's noise level is set as this base value multiplied by $(1 + c_v)$, where $c_v$ is
the coefficient of variation of the repetition time across that cohort's resting-state
acquisition records, computed from metadata. This translates a qualitative description of
scanner heterogeneity into a quantitative perturbation magnitude without picking a noise
level by hand. Two profiles are evaluated: an \ac{ADNI}-like profile using both feature
noise and edge perturbation at the calibrated level (reflecting multi-vendor, multi-model
heterogeneity), and an \ac{OASIS}-like profile using feature noise at the calibrated level
and edge perturbation at half that level (single-vendor but spanning materially different
scanner generations, so structural drift is plausible but expected to be smaller).

The scope of this simulation is explicitly bounded. It was designed before external
cross-cohort connectivity matrices existed, so it estimates rather than measures
cross-cohort degradation. The Gaussian and edge-rewiring perturbation is a proxy for
acquisition heterogeneity rather than a physical acquisition model, and it does not capture
systematic bias, geometric distortion or motion artefacts specific to any scanner. The
repetition-time coefficient of variation may mix resting-state protocol variants within a
cohort even after filtering. The simulation is reported as an estimate of the direction and
approximate magnitude of the degradation.

\paragraph{Interpretability validation design.}
An independent validation of the interpretability instrumentation is fixed prior to any
run. The validation has three components: a permutation null on the region map's overlap
with the \ac{DMN}, computed from 1,000 random reassignments of which 46 nodes carry the
\ac{DMN} label (preserving the true count) and reported as a percentile rank of the
observed overlap; Spearman correlation of the region map across seeds, as a measure of
reproducibility; and both statistics computed separately per cohort, so that a map
tracking cohort identity rather than disease is identified as such.

Three design constraints are recorded as method constraints rather than as results. First,
stability is measured cross-seed rather than cross-fold because only the best fold's map is
persisted per run, giving four maps per rung rather than twenty; this reduces the power of
the stability claim. Second, the overlap statistic is restricted to the 46 default-mode
regions of the Schaefer-200 cortical atlas, which contains no subcortical regions; the
hippocampal network is not assessable with this atlas. Third, the null hypothesis is a network-label
spin test rather than a subject-label permutation: because \ac{DMN} membership is an
anatomical label, no subject permutation changes which nodes belong to the \ac{DMN}, so
the null instead shuffles the label assignment across the 200 nodes while preserving the
count of 46 \ac{DMN} nodes.
